\documentclass[10.5pt]{ctexart}
\usepackage[a4paper,top=2.15cm,bottom=2.35cm,left=2.3cm,right=2.3cm,headsep=0.5cm,footskip=1.0cm]{geometry}
\usepackage{amsmath}
\usepackage{unicode-math}
\usepackage{xcolor}
\usepackage{array,tabularx}
\usepackage{enumitem}
\usepackage{titlesec,titletoc}
\usepackage{tcolorbox}
\tcbuselibrary{skins,breakable}
\usepackage{fancyhdr}
\usepackage{microtype}
\usepackage{needspace}
\usepackage{multicol}
\usepackage[hidelinks]{hyperref}
\hypersetup{pdftitle={语音信息处理 · 期末复习题库},pdfsubject={依据 2026 复习课件编写的六套复习题},pdfkeywords={语音信息处理,复习,题库}}

% ---------- 颜色 ----------
\definecolor{ink}{RGB}{28,29,32}
\definecolor{accent}{RGB}{30,72,110}
\definecolor{accentlt}{RGB}{126,152,178}
\definecolor{muted}{RGB}{112,120,130}
\definecolor{rule}{RGB}{203,211,219}
\definecolor{ansbg}{RGB}{243,247,251}
\definecolor{optink}{RGB}{58,61,67}
\definecolor{coverbg}{RGB}{247,249,252}

% ---------- 字体 ----------
\setmainfont{TeX Gyre Termes}
\setsansfont{Inter}
\setCJKmainfont{Noto Serif CJK SC}[BoldFont=Noto Sans CJK SC Bold, ItalicFont=Noto Serif CJK SC]
\setCJKsansfont{Noto Sans CJK SC}[BoldFont=Noto Sans CJK SC Bold]
\setCJKfamilyfont{zhsong}{Noto Serif CJK SC}
\setmathfont{TeX Gyre Termes Math}
\setmonofont{DejaVu Sans Mono}[Scale=0.86]

\color{ink}
\linespread{1.42}
\setlength{\parindent}{0pt}
\setlength{\parskip}{0.28em}
\setlength{\emergencystretch}{3em}
\setcounter{secnumdepth}{0}
\renewcommand{\sectionmark}[1]{\markboth{#1}{}}

% ---------- 页眉页脚 ----------
\pagestyle{fancy}
\fancyhf{}
\renewcommand{\headrulewidth}{0.35pt}
\renewcommand{\footrulewidth}{0pt}
\fancyhead[L]{\sffamily\footnotesize\color{muted}语音信息处理\quad 期末复习题库}
\fancyhead[R]{\sffamily\footnotesize\color{muted}\leftmark}
\fancyfoot[C]{\sffamily\footnotesize\color{muted}\thepage}
\fancypagestyle{plain}{\fancyhf{}\renewcommand{\headrulewidth}{0pt}\fancyfoot[C]{\sffamily\footnotesize\color{muted}\thepage}}

% ---------- 章首 ----------
\newcommand{\chapstart}[4]{%
  \clearpage
  \phantomsection\addcontentsline{toc}{section}{第 #1 章\quad #2}%
  \markboth{第 #1 章\quad #2}{}%
  \begin{tcolorbox}[enhanced jigsaw,breakable=false,colback=coverbg,colframe=coverbg,boxrule=0pt,
    left=10pt,right=10pt,top=9pt,bottom=9pt,arc=0pt,outer arc=0pt,before skip=0pt,after skip=1.1em]
    {\sffamily\footnotesize\color{accentlt}第 #1 套\ \ \textbar\ \ 本章共 #4 题}\par\vspace{0.25em}
    {\sffamily\bfseries\color{accent}\fontsize{20}{24}\selectfont 第 #1 章\quad #2}\par\vspace{0.5em}
    {\color{accent}\rule{\linewidth}{1.1pt}}\par\vspace{1.2pt}
    {\color{rule}\rule{\linewidth}{0.35pt}}\par\vspace{0.5em}
    {\sffamily\footnotesize\color{muted}#3}
  \end{tcolorbox}%
  \vspace{0.2em}\par}

% ---------- 题型 ----------
\titleformat{\subsection}[block]
  {\normalfont\sffamily\bfseries\color{accent}\fontsize{12.5}{15}\selectfont}
  {\raisebox{-0.06em}{\color{accent}\rule{3pt}{0.92em}}\hspace{0.55em}}{0em}{}
\titlespacing*{\subsection}{0pt}{1.55em plus 0.2em minus 0.15em}{0.65em}

% ---------- 题目列表 ----------
\newenvironment{qlist}
  {\begin{enumerate}[label={\sffamily\bfseries\color{accent}\arabic*.},leftmargin=2.4em,labelsep=0.55em,itemsep=0.62em,topsep=0.15em,parsep=0.3em,partopsep=0pt]}
  {\end{enumerate}}

% ---------- 选项 ----------
\newcommand{\optlist}[4]{%
  \par\vspace{0.34em}%
  {\small\color{optink}%
  \begin{list}{}{\leftmargin=1.8em\itemindent=-1.8em\topsep=0pt\partopsep=0pt\parsep=0.1em\itemsep=0.17em}%
  \item #1\item #2\item #3\item #4%
  \end{list}}}
\newcommand{\optlisttwo}[4]{%
  \par\vspace{0.34em}%
  {\small\color{optink}%
  \noindent\begin{tabularx}{\dimexpr\linewidth-0.2em\relax}{@{}>{\raggedright\arraybackslash}p{0.47\linewidth}@{\hspace{0.9em}}>{\raggedright\arraybackslash}X@{}}
  #1 & #2\\[0.28em] #3 & #4
  \end{tabularx}}}

% ---------- 答案 ----------
\newtcolorbox{ansbox}{enhanced,breakable,colback=ansbg,colframe=ansbg,boxrule=0pt,
  borderline west={2.3pt}{0pt}{accent!55!white},
  arc=0pt,outer arc=0pt,left=7.5pt,right=7pt,top=4.6pt,bottom=4.6pt,
  before skip=0pt,after skip=0.62em,fontupper=\small,fontlower=\small}
\newcommand{\anslabel}{{\sffamily\bfseries\color{accent}答\ 案：}\hspace{0.35em}}
\newcommand{\ans}[1]{\par\vspace{2.15\baselineskip}\begin{ansbox}\anslabel #1\end{ansbox}}
\newcommand{\ansL}[1]{\par\vspace{3.6\baselineskip}\begin{ansbox}\anslabel #1\end{ansbox}}

% ---------- 填空 ----------
\newcommand{\blank}{\rule[-0.42ex]{2.55em}{0.45pt}}
\newcommand{\blankl}{\rule[-0.42ex]{5.1em}{0.45pt}}

% ---------- 目录 ----------
\renewcommand{\contentsname}{目\ \ 录}
\titlecontents{section}[0pt]
  {\addvspace{0.62em}\sffamily\bfseries\color{accent}\fontsize{10.5}{13}\selectfont}
  {\thecontentslabel\hspace{0.4em}}{}
  {\hfill\color{accent}\contentspage}
\titlecontents{subsection}[1.7em]
  {\addvspace{0.06em}\sffamily\footnotesize\color{muted}}
  {\thecontentslabel\hspace{0.4em}}{}
  {\titlerule*[0.45pc]{\color{rule}.}\color{muted}\contentspage}

\begin{document}
\thispagestyle{empty}

\vspace*{2.1cm}
{\color{accent}\rule{\linewidth}{2pt}}\par\vspace{3pt}
{\color{rule}\rule{\linewidth}{0.4pt}}\par

\vspace{2.4cm}
{\sffamily\color{accentlt}\fontsize{11}{14}\selectfont 2026 复习课件\ \ ·\ \ 六章全覆盖}\par
\vspace{0.7cm}
{\sffamily\bfseries\color{accent}\fontsize{34}{40}\selectfont 语音信息处理}\par
\vspace{0.35cm}
{\sffamily\bfseries\color{ink}\fontsize{24}{28}\selectfont 期末复习题库}\par
\vspace{0.9cm}
{\color{accent}\rule{6cm}{1.2pt}}\par

\vfill

\begin{tcolorbox}[enhanced jigsaw,colback=coverbg,colframe=coverbg,boxrule=0pt,arc=0pt,outer arc=0pt,
  left=13pt,right=13pt,top=11pt,bottom=11pt]
  \small
  \begin{tabularx}{\linewidth}{@{}>{\sffamily\bfseries\color{accent}}l@{\hspace{1.1em}}X@{}}
    范\ \ 围 & 课件中标注「掌握」「基本概念」的内容，以及课件蓝字重点标注的概念；标注「了解」的内容不进入题库。\\
    结\ \ 构 & 六套题，每章一套：选择题、填空题、判断题、简答题。\\
    答\ \ 案 & 位于每题下方，题后留有作答空行，蓝色区段内为答案与解析。\\
    用\ \ 法 & 先闭卷做完整套，再对照答案；错题回到该章「总览（思维导图）」复查。\\
  \end{tabularx}
\end{tcolorbox}

\vspace{0.8cm}
{\sffamily\footnotesize\color{muted}依据《语音信息处理》2026 复习课件（第 1--6 章）编写}\par
\vspace{1.6cm}

\newpage
\thispagestyle{plain}
{\sffamily\bfseries\color{accent}\fontsize{18}{22}\selectfont 目\ \ 录}\par
\vspace{0.4em}
{\color{rule}\rule{\linewidth}{0.4pt}}\par
\vspace{0.6em}
{\parskip=0pt\relax\makeatletter\@starttoc{toc}\makeatother}


\chapstart{1}{语音产生与感知机理}{选择题 10 · 填空题 10 · 判断题 6 · 简答题 5}{31}

\subsection{一、选择题}

\begin{qlist}

\item 关于「语音」的定义，下列表述正确的是：
\optlist{A. 语音就是声音，两者完全等价}{B. 语音是声音的一种，是人在言语交流中由发音器官产生、由听觉系统感知的语言载体}{C. 语音是书写符号的读音，与发音器官无关}{D. 语音只能由机器合成}
\ans{B。声音泛指一切振动产生的声波；语音只是其中由人发音器官产生、经听觉系统感知的那一类，它同时具有生理物理属性（载体）和社会属性（承载内容）。}

\item 声带在发声过程中的作用是：
\optlisttwo{A. 只提供能量}{B. 既是阀门，又是振动部件}{C. 只起滤波作用}{D. 只承担辐射作用}
\ans{B。肺提供气流（能量来源），声带位于喉内、两声带之间的开口叫声门，既可开合控制气流（阀门），又因张弛振动产生准周期脉冲（振动部件）。}

\item 一端封闭、一端开放的声管（如外耳道），其第 $n$ 次谐振频率为：
\optlisttwo{A. $nc/2l$}{B. $(2n-1)c/4l$}{C. $nc/4l$}{D. $(2n-1)c/2l$}
\ans{B。两端开放或两端封闭的声管、两端固定的弦用 $f_n=\dfrac{nc}{2l}$；一端闭一端开的声管用 $f_n=\dfrac{(2n-1)c}{4l}$。外耳道属后者，取 $l\approx 2.5$ cm、$c\approx 350$ m/s，得 $f_1\approx 3.5$ kHz、$f_2\approx 10$ kHz。}

\item 成年人外耳道的初次谐振频率大约位于：
\optlisttwo{A. $1$ kHz}{B. $3.5$ kHz}{C. $10$ kHz}{D. $20$ kHz}
\ans{B。初次谐振约 $3.5$ kHz，二次谐振约 $10$ kHz 前后（课件原问即此）。}

\item 浊音的激励源可等效为：
\optlisttwo{A. 白噪声}{B. 准周期脉冲信号}{C. 单频正弦波}{D. 单位冲激}
\ans{B。浊音是声带张弛振动产生的准周期气流脉冲；清音是气流高速通过收缩处形成湍流，等效为白噪声；爆破音是声道先关闭、气压升高后突然打开。}

\item 元音的声学特征主要取决于：
\optlisttwo{A. 基频 $F_0$ 与 $F_1$}{B. 共振峰 $F_1$ 与 $F_2$}{C. 基音周期}{D. 声道长度}
\ans{B。在声学上元音特征主要取决于共振峰频率 $F_1$、$F_2$；$F_1$ 大小对应张口度，$F_2$ 大小对应舌位前后（圆唇会进一步拉低 $F_2$）。}

\item 音素与音位最本质的区别是：
\optlist{A. 音素数量少、音位数量多}{B. 音素按物理和生理属性划分，音位按社会和语言属性划分}{C. 音素是辅音，音位是元音}{D. 音素可以书写，音位不能书写}
\ans{B。音素是生成上可区分的最小语音单元，区分标准是物理和生理属性；音位是听觉上可区分的最小语音单元，区分标准是社会和语言属性。若干相近的音素可归为一个音位。}

\item 耳蜗内毛细胞的主要功能是：
\optlisttwo{A. 放大弱声}{B. 抑制过强声}{C. 把机械振动转换为电信号}{D. 维持身体平衡}
\ans{C。内毛细胞把耳蜗内淋巴液的振动转换为电信号，经听觉神经传到脑干、丘脑与听觉皮层；外毛细胞负责放大低水平声音、对过强声音做能动抑制；平衡由前庭系统负责。}

\item 描述人耳主观感觉声音「有多响」的物理量，其单位是：
\optlisttwo{A. Mel}{B. phon}{C. Hz}{D. dB}
\ans{B。响度用方（phon）表示，是主观心理量，与频率有关（等响度曲线）；音调的主观单位是梅尔（Mel）；Hz 与 dB 是客观量。}

\item 在基底膜上，高频与低频分别在何处引起最大振动：
\optlisttwo{A. 高频在蜗顶、低频在基部}{B. 高频在基部、低频在蜗顶}{C. 两者都在基部}{D. 与频率无关}
\ans{B。基底膜波动自耳蜗基部向蜗顶移动，不同频率在最匹配的一段振幅最大：高频靠基部、低频靠蜗顶，从而把频率转换成位置，耳蜗等效为一个滤波器组。}

\end{qlist}

\subsection{二、填空题}

\begin{qlist}

\item 语音是人在言语交流过程中由 \blank 产生、由 \blank 感知的语言载体。
\ans{发音器官；听觉系统。}

\item 声道作为 \blank ：放大某一频率而衰减其他频率分量；被放大的频率称为 \blank ／共振频率（formant frequency）。
\ans{谐振腔（谐振）；谐振频率。当激励的频率达到声道的固有频率，声道会以最大振幅振荡（共鸣）；共振峰由每一瞬间的声道外形决定，是声道的重要声学特征。}

\item \blank 是声带开启--闭合一次的时间，其倒数称为 \blank 。
\ans{基音周期（Pitch Period）$T_0$；基频（Pitch Frequency）$F_0=1/T_0$。（p26 蓝字：振幅 Amplitude、基音周期、基频、音调同为语音的基本特性。）}

\item 语音识别的建模单元可取 \blank 、\blank 、\blank 等。
\ans{音素；音节；字。（p40 蓝字；音节是语言的最基本单位，至少包括一个元音。）}

\item 人类语音的发音运动被 \blank 和符号化，这是同一语音单元在发音姿势变动时声学特性仍保持稳定的原因；语音识别的建模单元可取 \blank 、\blank 、\blank 等。
\ans{范畴化；音素；音节；字。（「发音运动指令的范畴化」与「语音识别的建模单元」均为课件蓝字。）}

\item 中耳鼓膜振动经 \blank 传递到耳蜗的 \blank ；\blank 对 $500\sim 2000$ Hz 的声波产生较大的谐振，起到带通滤波器的作用。
\ans{听骨链（锤骨 $\to$ 砧骨 $\to$ 镫骨）；前庭窗；听骨链。}

\item 外耳道的初次谐振频率大约位于 \blank 附近，二次谐振频率将达 \blank 前后。
\ans{3.5 kHz；10 kHz。（外耳道长约 2.5 cm、常温下音速约 350 m/s，按一端闭一端开声管公式 $f_n=\dfrac{(2n-1)c}{4l}$ 计算。）}

\item 耳蜗由三个充满淋巴液的腔组成：\blank 、\blank 、\blank ；\blank 把耳蜗内液体的声音振动转化为电学信号，\blank 则机械地放大传入的低水平声音、并对过强声音进行能动的抑制。
\ans{前庭阶、蜗管、鼓阶；内毛细胞；外毛细胞。（p47 蓝字。）}

\item \blank ：声音小到人耳刚刚能听见时的大小；\blank ：声音大到人耳刚刚感觉痛时的大小。
\ans{听阈；痛阈。响度反映一个人主观感觉不同频率成分的声音强弱的物理量，是主观心理量、与频率有关。}

\item 言语链的 \blank ：\blank ；\blank ：把机械振动转换为神经脉冲，经听觉神经传到听觉脑干、丘脑和听觉皮层。
\ans{产生——言语（语音）产生；感知——言语（语音）感知；言语链反向运作。（言语链反映说话人意图如何编码成言语传给听话人、听话人如何解码再现意图。）}

\end{qlist}

\subsection{三、判断题}

\begin{qlist}

\item 语音与声音是同一概念，凡是声音都是语音。
\ans{错误。语音只是声音的一种，须由人的发音器官产生、由听觉系统感知，并承载语言信息。}

\item 声道的谐振频率由每一瞬间的声道外形决定，因此调整声道形状即可发出不同的音。
\ans{正确。调音就是为发出不同声音而调整声道形状；共振频率（共振峰）随声道外形改变。}

\item 辅音一般没有共振峰，只有半元音可以被认为有共振峰。
\ans{正确。这是课件明确给出的结论。}

\item 一个音位只能对应一个音素。
\ans{错误。关系是「若干相近的音素归为一个音位」，即一个音位可以包含多个音素。}

\item 高频声波在基底膜上靠近蜗顶处引起最大振动。
\ans{错误。高频靠基部、低频靠蜗顶。}

\item 痛阈是指声音小到人耳刚刚能听见的声压级。
\ans{错误。那是听阈；痛阈是声音大到人耳刚刚感觉痛。}

\end{qlist}

\subsection{四、简答题}

\begin{qlist}

\item 简述语音产生的两个基本步骤，并说明「声源」与「声道」各自的作用。
\ansL{两个步骤：①声门／声带产生不同频率的声音（准周期气流脉冲或白噪声）；②声道对声源进行调制，再经唇、鼻孔的声波放射与声道壁振动辐射形成语音。声源决定「怎么响」：浊音为声带张弛振动形成的准周期脉冲，清音为湍流形成的白噪声，爆破音为声道突然打开形成的瞬态；声道作为谐振腔决定「响成哪个音」：它放大某些频率、衰减其他频率，被放大的频率即共振峰。}

\item 音素与音位有何区别？请从定义、区分标准和数量关系三方面作答。
\ansL{定义：音素是生成上可区分的最小语音单元，音位是听觉上可区分的最小语音单元。区分标准：音素依据物理和生理属性，音位依据社会和语言属性（约定俗成）。数量关系：若干相近的音素归类为一个音位，故音位数不多于音素数；全世界约 200 个元音、600 个辅音，而一种特定语言大约只有 40 个音素。例：普通话中 z、c、s 与 zh、ch、sh 是两组不同音位，而在许多南方方言中属同一组音位。}

\item 判断元音与辅音的依据是什么？它与「浊音／清音」的划分是同一套标准吗？
\ansL{依据是气流在声道中是否受阻：元音由相对开放的声道产生，有声带振动但没有可听见的摩擦音；辅音是发音器官形变使气流在声道中部分或完全受阻而产生的语音。它与「浊／清」不是同一套划分：浊／清看声带是否振动，元／辅看声道是否受阻，两者相互独立，例如浊辅音同时满足「声带振动」与「声道受阻」。}

\item 写出言语链中「言语产生」方向的四个环节，并说明它与哪些人机语音交互任务相对应。
\ansL{（课件标「掌握」）产生方向四环节：①基于意图决定说话内容，言语中枢进行言语规划（把抽象概念映射到词汇表征，并给出音素、语调和时长的正确顺序）；②大脑言语运动中心编排发音运动程序，按时序传到下一级神经系统；③运动指令传到发声肌肉，肌肉收缩使肺部气流通过声带，产生不同类型的声源；④声源经声道调制，再由唇、鼻孔的声波放射和声道壁振动辐射产生语音信号。对应关系：人的语音感知对应 ASR（语音识别），人的语音产生对应 TTS（语音合成），辨识说话人对应声纹识别（说话人识别）。感知方向则反向运作：机械振动 $\to$ 流体行波 $\to$ 神经脉冲 $\to$ 听觉皮层 $\to$ 语言中枢解码音素、再现意图。}

\item 外耳、中耳、内耳各承担什么功能？为什么说内耳具有双重功能？
\ansL{外耳（耳廓＋外耳道，长约 2.5 cm）负责集声，耳廓使频谱峰压点 5.5 kHz 的纯音提高约 10 dB，外耳道作为一端闭一端开的声管产生 3.5 kHz 与 10 kHz 的谐振。中耳（鼓膜＋听骨链＋咽鼓管）放大来自外耳的振动，经前庭窗传入内耳，对 500\textasciitilde 2000 Hz 谐振，起带通滤波器作用。内耳骨迷路包括耳蜗与前庭系统：前庭系统是平衡觉的末梢器官，耳蜗把来自外耳和中耳的机械振动转换为神经电信号传给大脑，故内耳兼有听觉与感受位置变化的双重功能。}

\end{qlist}


\chapstart{2}{语音信号处理的基础}{选择题 11 · 填空题 10 · 判断题 6 · 简答题 5}{32}

\subsection{一、选择题}

\begin{qlist}

\item 奈奎斯特频率是指：
\optlisttwo{A. 人耳能听到的最高频率}{B. 离散信号系统采样频率的一半}{C. 语音信号的基频}{D. 滤波器的截止频率}
\ans{B。奈奎斯特频率为 $f_s/2$。奈奎斯特--香农采样定理指出：给定采样频率 $f_s$，能完全重构的频带限制为 $B\le f_s/2$，只要奈奎斯特频率高于被采样信号的最高频率（带宽），就不会发生混叠。}

\item 语音信号预处理中加窗的主要目的是：
\optlist{A. 提高信噪比}{B. 满足离散傅里叶变换的周期性要求、减少泄漏}{C. 去除噪声}{D. 压缩数据量}
\ans{B。加窗后信号 $x'(n)=x(n)w(n)$，使时域信号更好地满足 DFT 的周期性要求、减少泄漏。常用汉明窗、汉宁窗、矩形窗（矩形窗旁瓣泄漏大）。}

\item 分帧时窗长一般取 $20\sim 25$ ms，主要原因是：
\optlisttwo{A. 计算方便}{B. 约等于 $1\sim 7$ 个基音周期}{C. 与采样率成正比}{D. 人耳时间分辨率上限}
\ans{B。人的基音范围 $80\sim 500$ Hz，对应基音周期 $2\sim 12.5$ ms；窗长取 $1\sim 7$ 个基音周期，既能包含足够多的周期，又不至于跨越音素变化。（课件原问「为什么分帧的窗长一般取 25 ms」即此。）}

\item 短时能量的表达式 $E_n=\sum_m x^2(m)h(n-m)$ 中，$h(n)$ 等于：
\optlisttwo{A. $w(n)$}{B. $w^2(n)$}{C. $|w(n)|$}{D. $\log w(n)$}
\ans{B。因为 $E_n=\sum_m [x(m)w(n-m)]^2=\sum_m x^2(m)w^2(n-m)$。短时能量可用于浊音/清音区分、端点检测、声母与韵母分界、情感分类特征。}

\item 短时平均过零率（ZCR）的典型值为：
\optlist{A. 浊音约 4500 Hz、清音约 1400 Hz}{B. 浊音约 1400 Hz、清音约 4500 Hz}{C. 两者均约 3000 Hz}{D. 与浊清无关}
\ans{B。ZCR 反映包括主能量在内的区域的平均频率，浊音约 1400 Hz、清音约 4500 Hz；常与短时能量结合用于端点检测。}

\item 关于短时自相关函数，下列说法错误的是：
\optlist{A. 具有对称性 $R(m)=R(-m)$}{B. $m=0$ 时取最大值}{C. 白噪声的自相关只在 $m=0$ 处有强峰}{D. 周期信号的自相关在所有 $m$ 上恒定}
\ans{D。若 $x(n)$ 是周期信号，$R(i)$ 会在 $i$ 等于周期整数倍处取极大值，这正是用它估计基音周期的依据。}

\item 关于语谱图，下列说法正确的是：
\optlist{A. 宽带语谱图（窗长约 5 ms）呈现水平条纹}{B. 窄带语谱图（窗长约 40 ms）呈现垂直条纹}{C. 宽带语谱图时间分辨率高，垂直条纹间隔为基音周期，适合看共振峰}{D. 窗长越长，时间分辨率越高}
\ans{C。宽带（约 200 Hz／5 ms）频率分辨率低、时间分辨率高，呈垂直条纹，间隔即基音周期，用于观测共振峰；窄带（约 25 Hz／40 ms）频率分辨率高，呈水平条纹，用于观测谐波与基音频率。核心权衡是时间分辨率与频率分辨率不能同时高。}

\item 倒谱的定义是：
\optlist{A. 信号的傅里叶变换}{B. 信号傅里叶变换经对数运算后再作傅里叶反变换得到的谱}{C. 信号的自相关函数}{D. 信号的功率谱}
\ans{B。语音中一般先取频谱幅度。目的是把声源与声道的卷积（频域相乘）经对数变成相加，从而在倒谱域把它们分开。}

\item 在倒谱中，$C(n)\ge F_0$ 与 $C(n)<F_0$ 分别主要对应：
\optlisttwo{A. 声道信息与声源信息}{B. 声源信息与声道（系统）信息}{C. 噪声与语音}{D. 低频与高频}
\ans{B。声源频谱变化比声道频谱快得多，因此低倒谱段反映声道包络、高倒谱段反映声源（基音）；用 Lifter 窗取出低倒谱段即得声道包络。}

\item 线性预测编码（LPC）将声道建模为：
\optlisttwo{A. 全零点系统}{B. 全极点（自回归 AR）系统}{C. 带通滤波器组}{D. 有限脉冲响应系统}
\ans{B。$H(z)=\dfrac{G}{1-\sum_{i=1}^{p}a_iz^{-i}}$，$G$ 为增益、$a_i$ 为实数系数、$p$ 为阶数；$a_i$ 通过最小化残差均方误差 $E[\varepsilon^2(n)]$ 求得。}

\item 梅尔频率倒谱系数（MFCC）提取流程中，第 5 步是：
\optlisttwo{A. 取对数}{B. 离散余弦变换（DCT）}{C. Mel 滤波}{D. 预加重}
\ans{B。六步为：预加重、分帧、加窗 $\to$ DFT 取模平方得能量谱 $\to$ Mel 三角滤波器组 $\to$ 取对数 $\to$ DCT $\to$ 取低阶系数。其中 $C_0$ 表示语音的平均强度，语音识别中一般不直接使用。}

\end{qlist}

\subsection{二、填空题}

\begin{qlist}

\item 常用的采样频率有 \blank kHz、\blank kHz、\blank kHz，其中 CD 音质使用 \blank kHz。
\ans{8；16；44.1；44.1。原因：人的听觉范围为 $20\sim 20000$ Hz，$44.1$ kHz 的奈奎斯特频率为 $22.05$ kHz，足以覆盖。}

\item 如果 $x(n)$ 是一个周期信号，$R(i)$ 就会在 \blank 处取极大值——这是短时自相关函数可以应用到周期估计等任务的原因。
\ans{$i$ 等于周期整数倍的点。（p23 蓝字。）}

\item \blank （zero-cross rate，ZCR）：信号采样点符号变化的次数；\blank （Autocorrelation）：测定相同信号在时域内的相似度。
\ans{短时间平均过零率；自相关函数。短时自相关函数在自相关函数的基础上加短时分析。}

\item 已知系统的输入和脉冲响应可求输出；已知输入和输出可辨识系统；语音信号处理多属于第 \blank 种情况。
\ans{3（只有输出已知）。}

\item \blank ：信号的傅里叶变换经对数运算后再进行傅里叶反变换得到的谱（语音中一般对频谱取幅度值）。
\ans{倒谱（Cepstrum）。}

\item 用系数 \blank 可以定义一个 \blank 阶 \blank （Linear Prediction，线性预测）。
\ans{$\{a_i\}$；$p$；线性预测器。（p57 蓝字：将语音 $x(n)$ 看作声源 $s(n)$ 激励全极点系统 $H(z)$ 的输出。）}

\item \blank （带通滤波器）实现听觉末梢系统中的声学频率分析；\blank 是根据心理物理实验结果获得的滤波器带宽。
\ans{听觉滤波器；等效矩形带宽（ERB, Equivalent Rectangular Bandwidth）。}

\item \blank （ITD）：来自右侧的声音到达右耳比左耳早；\blank （IID）：头部遮挡使远侧耳强度更低。
\ans{两耳时间差（Interaural Time Difference）；两耳声强差（Interaural Intensity Difference）。}

\item 语音信号的频谱分析，是 \blank 和 \blank 语音信号的重要方法。
\ans{认识；处理。（p25 蓝字原句。）}

\item \blank （200 Hz，窗长 5 ms）频率分辨率低、时间分辨率高；\blank 频率分辨率高、时间分辨率低。
\ans{宽带滤波器；窄带滤波器。}

\end{qlist}

\subsection{三、判断题}

\begin{qlist}

\item 只要采样频率足够高，任何信号都可以被完全重构而不产生混叠。
\ans{正确（在采样定理前提下）。条件是信号带宽 $B\le f_s/2$，即奈奎斯特频率高于信号最高频率或带宽。}

\item 矩形窗的旁瓣泄漏比汉明窗小，因此更适合做语音分帧。
\ans{错误。矩形窗旁瓣泄漏大，汉明窗、汉宁窗可压低旁瓣。}

\item 短时能量和短时平均过零率常常结合使用来完成语音端点检测。
\ans{正确。能量反映强弱，ZCR 反映频率高低，两者互补。}

\item 语谱图的窗长越长，频率分辨率越高、时间分辨率也越高。
\ans{错误。窗越长频率分辨率越高，但时间分辨率下降，两者不能同时提高。}

\item 倒谱分析中取对数的目的是把频域的相乘关系转换为相加关系。
\ans{正确。$X(k)=S(k)H(k)\Rightarrow \log X(k)=\log S(k)+\log H(k)$，从而可用 Lifter 分离声源与声道。}

\item LPC 与倒谱都直接利用声源--滤波器模型求解声道参数。
\ans{错误。LPC 直接利用声源--滤波器模型得到声道滤波器系数；倒谱是间接地通过时频分析把声源与声道分开。}

\end{qlist}

\subsection{四、简答题}

\begin{qlist}

\item 为什么要对语音信号分帧？为什么帧长取 $20\sim 25$ ms、并且相邻帧之间要有帧移？
\ansL{语音是典型的非周期信号，整体特性随时间变化，但在 $20\sim 30$ ms 的短时间内在发音器官惯性作用下近似准周期、特性稳定，因此可以切成短帧分别处理，并可近似为线性时不变系统。帧长的取法依据基音周期：人基音范围 $80\sim 500$ Hz，周期 $2\sim 12.5$ ms，窗长取 $1\sim 7$ 个基音周期，既能包含足够周期数，又不跨越音素变化。设置帧移使相邻帧重叠，避免帧边界处信息丢失，保证特征随时间的连续性。}

\item 简述倒谱分析的原理，并说明为什么它能分离声源与声道。
\ansL{原理：先把语音看成为声源与声道冲激响应的卷积 $x(n)=\sum_i s(i)h(n-i)$，频域上变为乘积 $X(k)=S(k)H(k)$；取对数得 $\log X(k)=\log S(k)+\log H(k)$，把相乘变为相加；再作傅里叶反变换（IFFT）得到倒谱 $C(n)$，落在类频率（quefrency）域。分离依据：声源频谱随频率变化快、声道频谱变化慢，因此在倒谱域中低倒谱段对应声道包络、高倒谱段对应声源基音，以 $F_0$ 为分界；用 Lifter 取出低倒谱段再作傅里叶变换即得声道包络。流程为：语音 $\to$ 预处理 $\to$ 傅里叶变换 $\to$ 取对数 $\to$ 逆傅里叶变换 $\to$ 去高阶倒谱 $\to$ 傅里叶变换 $\to$ 包络。}

\item 写出 MFCC 的提取步骤，并说明 $C_0$ 的含义。
\ansL{六步：①预处理——预加重（提升高频，补偿辐射造成的每倍频程约 6 dB 跌落）、分帧、加窗；②对每帧作 DFT，取模的平方得到离散能量谱，记 $S(k)=|X(k)|^2$；③通过一组 Mel 三角滤波器组，得到各 Mel 带的能量 $m(i)$；④取对数 $\lg[m(i)]$；⑤作离散余弦变换（DCT），得 $C_{\mathrm{MFCC}}(n)=\sum_i \lg[m(i)]\cos\frac{n(i-0.5)\pi}{L}$；⑥取低阶系数作为特征向量。Mel 刻度用 $\mathrm{Mel}(f)=2595\log_{10}(1+f/700)$ 模拟人耳非线性感知。$C_0$ 表示语音的平均强度，语音识别中一般不直接使用。}

\item 简述线性预测编码（LPC）的基本思路，并说明残差与系数 $a_i$ 的物理含义。
\ansL{基本思路：把语音 $x(n)$ 看作声源 $s(n)$ 激励一个全极点（AR）系统 $H(z)=\dfrac{G}{1-\sum_{i=1}^{p}a_iz^{-i}}$ 的输出，用前 $p$ 个采样线性组合预测当前值 $\hat x(n)=\sum_{i=1}^{p}a_ix(n-i)$。残差 $\varepsilon(n)=x(n)-\hat x(n)$ 是「预测不出来的那部分」，对应声源激励；系数 $a_i$ 刻画声道滤波器的极点位置，即声道的「指纹」。最优 $a_i$ 通过最小化残差均方误差 $E[\varepsilon^2(n)]$ 求得。}

\item 比较宽带与窄带语谱图，并说明时间分辨率与频率分辨率的权衡关系。
\ansL{宽带语谱图带宽约 200 Hz、窗长约 5 ms：频率分辨率低、时间分辨率高，呈垂直条纹且间隔等于基音周期，适合观察共振峰与浊音的周期性。窄带语谱图带宽约 25 Hz、窗长约 40 ms：频率分辨率高、时间分辨率低，呈水平条纹，适合观察谐波 $H_1,H_2,\dots$ 与基音频率。权衡关系：$B\cdot T$ 近似为常数，窗越长频率越准但时间越糊，二者不能同时提高，实际按任务需要折中选取窗长；这也是短时傅里叶变换的固有局限。}

\end{qlist}


\chapstart{3}{语音识别}{选择题 14 · 填空题 12 · 判断题 6 · 简答题 6}{38}

\subsection{一、选择题}

\begin{qlist}

\item 自动语音识别（ASR）的任务是：
\optlisttwo{A. 由文本生成语音}{B. 把语音信号转变为相应的文本}{C. 判断说话人身份}{D. 去除背景噪声}
\ans{B。识别结果可作最终输出（转写），也可作为自然语言处理（对话、翻译）的输入。}

\item 按发音方式划分，语音识别可分为：
\optlist{A. 小、中、大词汇表识别}{B. 特定说话人与非特定说话人识别}{C. 孤立词、连接词、连续语音识别}{D. 单通道与多通道识别}
\ans{C。三个分类维度是：发音方式（孤立词／连接词／连续语音）、词汇表大小（10\textasciitilde 100／100\textasciitilde 1000／1000 词条以上）、说话人（特定／非特定）。}

\item 词错误率（WER）的分子由哪三类错误构成：
\optlisttwo{A. 插入、替换、删除}{B. 漏检、误检、重复}{C. 训练、解码、对齐}{D. 声学、语言、词典}
\ans{A。$\mathrm{WER}=\dfrac{100\times(\#I+\#S+\#D)}{\text{正确转写中的总词数}}$，其中 Insertion 为插入、Substitution 为置换、Deletion 为删除。}

\item 动态时间弯折（DTW）的核心思想是：
\optlist{A. 用神经网络分类}{B. 用动态规划把两个不等长的矢量序列对齐后再算相似度}{C. 用统计模型估计概率}{D. 用语言模型约束结果}
\ans{B。DTW 建立从 $(1,1)$ 到 $(I,J)$ 的路径，累积距离最小的路径即最优对齐路径，再用路径长度归正。其缺点：难区分相似模板、声学条件不同时性能急剧下降、未充分利用时序动态信息、不利于大词汇表连续语音识别。}

\item 噪音通道模型经贝叶斯法则化简后的识别判据是：
\optlisttwo{A. $\arg\max_W P(W\mid X)=\arg\max_W P(X\mid W)P(W)$}{B. $\arg\max_W P(X)$}{C. $\arg\max_W P(W)$}{D. $\arg\max_W P(X\mid W)/P(W)$}
\ans{A。$\arg\max_W\dfrac{P(X\mid W)P(W)}{P(X)}$，因 $P(X)$ 与 $W$ 无关可去掉分母。$P(X\mid W)$ 是声学模型、$P(W)$ 是语言模型。（课件标「掌握」，框架图标注「重点掌握」。）}

\item 隐马尔可夫模型的「隐」是指：
\optlist{A. 观测值不可见}{B. 状态序列不可直接观测，只能通过观测序列推断}{C. 参数不可估计}{D. 转移概率未知}
\ans{B。HMM 用双重随机过程描述：一个 Markov 链描述状态转移（不可见），另一个随机过程描述状态与观察值的统计对应关系（可见）。}

\item HMM 的三个基本问题与算法的对应关系是：
\optlist{A. 概率计算—Viterbi，解码—前向算法，参数估计—Baum-Welch}{B. 概率计算—前向／后向算法，解码—Viterbi 算法，参数估计—Baum-Welch 算法}{C. 三者都用 Viterbi}{D. 三者都用 EM}
\ans{B。问题 1 计算 $P(O\mid\lambda)$ 用前向／后向算法；问题 2 求最佳状态序列 $Q^*$ 用 Viterbi；问题 3 调整 $\lambda$ 使 $P(O\mid\lambda)$ 最大用 Baum-Welch。（课件标「掌握」。）}

\item 前向算法的计算量是：
\optlisttwo{A. $O(NT)$}{B. $O(N^2T)$}{C. $O(N^3T)$}{D. $O(2^TN)$}
\ans{B。当 $N=5$、$T=100$ 时只需约 2500 次乘法，而直接计算法需 $O(N^T)$，这是引入前向／后向算法的根本原因。}

\item Viterbi 算法与前向算法的根本区别在于：
\optlist{A. Viterbi 求和、前向取最大}{B. Viterbi 每步取最大并记录回溯指针 $\varphi$，前向对所有路径求和}{C. 两者完全相同}{D. Viterbi 只能用于离散 HMM}
\ans{B。$\delta_t(j)=\max_i[\delta_{t-1}(i)a_{ij}]b_j(o_t)$，$\varphi_t(j)=\arg\max_i[\delta_{t-1}(i)a_{ij}]$；终了取 $\max_i\delta_T(i)$ 后按 $\varphi$ 回溯得 $Q^*$。}

\item Baum-Welch 算法的性质是：
\optlist{A. 能保证求得全局最优参数}{B. 是期望最大化（EM）算法的一个特例，未必收敛到全局最优}{C. 用于解码}{D. 用于计算观测概率}
\ans{B。它用于观测序列与隐状态序列都未知时的参数估计，通过 E 步与 M 步迭代提高 $P(O\mid\lambda)$，可能收敛到局部极值。}

\item 大词汇表连续语音识别（LVCSR）系统的三层框架是：
\optlisttwo{A. 声学层、韵律层、语义层}{B. 声学—语音层、词层、句法层}{C. 前端、后端、解码}{D. 特征层、模型层、决策层}
\ans{B。声学—语音层以子词（subword）为单位输出，每个子词对应一套 HMM 结构与参数；词层规定每个词由哪些音素／音子串接而成（词典）；句法层规定词按什么规则组合成句子（语言模型）。}

\item 三音素建模的记法与参数量级是：
\optlisttwo{A. $c\text{-}l\text{+}r$，$N^2$}{B. $l\text{-}c\text{+}r$，$N^3$}{C. $l\text{+}c\text{-}r$，$N$}{D. $l\text{-}c\text{+}r$，$N^2$}
\ans{B。$c$ 为中心单元、$l$ 为左相关、$r$ 为右相关信息；$N$ 个音素共需 $N^3$ 个三音素，参数量过大，故实际用决策树聚类共享参数，该过程依赖单音素建模后的对齐。}

\item 与 GMM 相比，DNN-HMM 中输入特征的特点是：
\optlist{A. 只用单帧 MFCC}{B. 多采用滤波器组（FBK）特征并拼接相邻多帧}{C. 只用时域波形}{D. 必须做矢量量化}
\ans{B。DNN 会拼接相邻若干帧以利用更多上下文；输出维度对应 HMM 状态数（通常每维对应一个绑定的 triphone 状态）；训练标注需先用 GMM-HMM 做强制对齐（Viterbi）给每帧打状态标签。}

\item 语言模型的困惑度（perplexity）：
\optlist{A. 越高越好}{B. 越低越好，等于模型分配给测试集概率的 $-1/N$ 次方}{C. 与句长无关}{D. 等于 WER}
\ans{B。$PP(W)=P(w_1w_2\cdots w_N)^{-1/N}$，含义是预测下一个词时的加权平均分支因子，越低越好。例（WSJ 语料）：Unigram 962、Bigram 170、Trigram 109。}

\end{qlist}

\subsection{二、填空题}

\begin{qlist}

\item HMM 的参数组可写为 $\lambda=(\pi,A,B)$，其中 $\pi$ 为 \blank ，$A$ 为 \blank ，$B$ 为 \blank 。
\ans{初始状态概率分布；状态转移概率矩阵；观测值概率矩阵。（完整写法为 $\lambda=(S,V,\pi,A,B)$。）}

\item \blank ：到时刻 $t$ 为止部分观测序列为 $o_1\cdots o_t$ 且状态为 $s_i$ 的概率。
\ans{前向变量，$\alpha_t(i)=P(o_1\cdots o_t,\ q_t=i\mid\lambda)$。}

\item \blank 定义为 $P(o_{t+1}\cdots o_T\mid q_t=i,\lambda)$，对应的算法称为 \blank 算法。
\ans{后向变量；后向。}

\item 前向／后向算法的计算量为 $O($ \blank $)$；当 $N=5$、$T=100$ 时，只需约 \blank 次乘法运算。
\ans{$N^2T$；2500。}

\item 后向变量定义为 \blank （用语言叙述）；其算法计算量为 \blank 数量级。
\ans{$P(o_{t+1}\cdots o_T\mid q_t=i,\lambda)$（从 $t+1$ 到 $T$ 的观测且 $t$ 时刻状态为 $i$ 的概率）；$N^2T$。}

\item 语音识别解码可采用 \blank 搜索算法；Viterbi 算法最终得到最佳状态序列 \blank 以及所对应的概率 \blank 。
\ans{Viterbi；$Q^*$；$P^*$。}

\item 高斯混合模型可以看作是由 \blank 个单高斯模型组合而成的模型，用来描述复杂的 \blank 分布。
\ans{$M$；数据（如语音数据）。}


\item N-gram 的近似假设是：词 $w_k$ 出现的概率只与前 \blank 个词有关，即 \blank 阶马尔可夫模型。
\ans{$N-1$；$N-1$。故 $P(w_1^n)=\prod_{k=1}^{n}P(w_k\mid w_{k-N+1}^{k-1})$。}

\item 困惑度定义为 $PP(W)=$ \blankl 。
\ans{$P(w_1w_2\cdots w_N)^{-1/N}=P(w_1^N)^{-1/N}$，越低越好。}

\item CTC 是 \blank 的缩写，即基于 \blank 的连接时序分类；传统声学模型训练必须预先训练一个 \blank 声学模型、采用其进行自动对齐标注。
\ans{Connectionist temporal classification；神经网络；GMM-HMM。（p146 蓝字。）}

\item \blank 是指在语流中，音段并非是静止的、分离的声音，音段会对相邻的音段产生影响；解决方法是 \blank 建模。
\ans{协同发音（coarticulation）；三音素（考虑音素的上下文，表示为 $l\text{-}c\text{+}r$，$c$ 为中心单元、$l$ 为左相关信息、$r$ 为右相关信息）。}

\end{qlist}

\subsection{三、判断题}

\begin{qlist}

\item 马尔可夫性质（无后效性）是指：在给定当前信息的条件下，过去的历史对预测将来无关。
\ans{正确。一阶马尔可夫链满足 $P(q_t\mid q_{t-1},q_{t-2},\dots)=P(q_t\mid q_{t-1})$。}

\item HMM 的前向算法与 Viterbi 算法的递推结构完全相同，只是记号不同。
\ans{错误。前向算法对所有可能的路径求和，得到 $P(O\mid\lambda)$；Viterbi 每一步只保留最优路径并记录回溯指针，得到最优状态序列。}

\item Baum-Welch 算法一定能求得 HMM 参数的全局最优解。
\ans{错误。它可能收敛到局部极值。}

\item 以词为建模单元比以音素为建模单元更节省存储与计算。
\ans{错误。以词为单元会造成大量冗余存储和计算；一般采用比词小的子词（音节、半音节、音素）作为识别基元。}

\item 声学单元越小，对上下文的敏感性越大，越容易受相邻音影响而产生变异。
\ans{正确。单元小则数量少、训练工作量小，但对上下文更敏感、类型设计与样本采集更困难。}

\item 困惑度越低说明语言模型对测试数据的拟合越好。
\ans{正确。困惑度是预测下一个词时的加权平均分支因子，越低越好。}

\end{qlist}

\subsection{四、简答题}

\begin{qlist}

\item 写出 HMM 的三个基本问题，并指出各自的求解算法。
\ansL{问题 1（模型概率计算）：已知模型参数 $\lambda=(\pi,A,B)$ 与观测序列 $O=o_1o_2\cdots o_T$，计算 $P(O\mid\lambda)$。解：前向算法（递推 $\alpha_t(j)=[\sum_i\alpha_{t-1}(i)a_{ij}]b_j(o_t)$）或后向算法（$\beta_t(i)=\sum_j\beta_{t+1}(j)a_{ij}b_j(o_{t+1})$），计算量均为 $O(N^2T)$。问题 2（模型解码）：给定 $\lambda$ 与 $O$，求最佳状态序列 $Q^*=\arg\max_Q P(O,Q\mid\lambda)$。解：Viterbi 算法。问题 3（模型参数估计）：给定观测序列集合，调整参数 $\lambda$ 使 $P(O\mid\lambda)$ 最大。解：Baum-Welch 算法（EM 的特例）。}

\item 为什么前向算法把计算量从 $O(N^T)$ 降到 $O(N^2T)$？它与后向算法、Viterbi 算法分别解决什么问题？
\ansL{直接计算法需要枚举所有 $N^T$ 条长度为 $T$ 的状态路径并求和，计算量随 $T$ 指数增长，不可行。前向算法利用动态规划思想：把「到 $t$ 时刻观测为 $o_1\dots o_t$ 且状态为 $i$」的概率作为中间量 $\alpha_t(i)$ 递推，每步只需对 $N$ 个前驱状态求和，故总复杂度为 $O(N^2T)$（如 $N=5$、$T=100$ 只需约 2500 次乘法）。三者的分工：前向算法从 $t=1$ 向后算出 $P(O\mid\lambda)$；后向算法从 $t=T$ 向前算出同样的量，两者常在 Baum-Welch 中配合使用；Viterbi 算法同样逐帧推进，但每步取最大而非求和，用于求最优状态序列（解码）。}

\item 为什么用高斯混合模型（GMM）作为 HMM 的观测值概率分布？它相比离散 HMM 有何优势？
\ansL{离散 HMM 需先做矢量量化（VQ），把连续特征映射到有限的码字，量化必然引入误差且难以刻画复杂分布。GMM 用 $M$ 个高斯分布的加权和 $P(x\mid\theta)=\sum_m c_m\mathcal N(x;\mu_m,\Sigma_m)$ 直接对连续特征建模：高斯分布具有优良的数学性质与计算性能，加权和可以平滑地逼近任意形状的概率密度函数，对实际数据具备较强的表征能力。放进 HMM 的某个状态 $i$ 即得 $b_i(o_t)=\sum_m c_{i,m}\mathcal N(o_t;\mu_{i,m},\Sigma_{i,m})$，其中 GMM 负责「这一帧像哪个状态」、HMM 负责「这些状态按什么顺序出现」。}

\item 建模单元如何选择？为什么需要三音素建模，又为什么要用决策树？
\ansL{选择原则：以词为单元不合理（冗余存储与计算过大），因此采用比词小的子词（音节、半音节、音素等）。权衡：声学单元越小，数量越少、训练工作量越小；但单元越小对上下文越敏感，越容易受前后相邻音影响产生变异，类型设计与样本采集更困难。单音素建模没有考虑协同发音（语流中音段会互相影响），解决办法是用三音素 $l\text{-}c\text{+}r$ 把左右上下文写进模型。$N$ 个音素需要 $N^3$ 个三音素，参数量随 $N$ 急剧增大，因此用决策树对三音素的高斯分布做聚类以共享参数、缓解数据稀疏；该过程依赖单音素建模后的对齐。从单音素到三音素的流程为：单音素＋EM 训练 $\to$ 把高斯分布复制到三音素 $\to$ 建决策树并聚类 $\to$ 复制并混合训练。}

\item 简述 N-gram 语言模型的基本假设与概率估计方法，并说明困惑度的含义。
\ansL{链式法则给出 $P(w_1^n)=\prod_{k=1}^{n}P(w_k\mid w_1^{k-1})$，但完整上下文所需参数量随上下文长度指数增长，无法估计。N-gram 近似假设词 $w_k$ 只与前 $N-1$ 个词有关（$N-1$ 阶马尔可夫模型）：$P(w_1^n)=\prod_{k=1}^{n}P(w_k\mid w_{k-N+1}^{k-1})$；Unigram 为 $P(\text{phone})$、Bigram 为 $P(\text{phone}\mid\text{cell})$、Trigram 为 $P(\text{phone}\mid\text{your cell})$。参数用相对频率估计：$P(w_n\mid w_{n-1})=C(w_{n-1}w_n)/C(w_{n-1})$；为使概率模型一致，句子两端需加开始符与结束符。困惑度 $PP(W)=P(w_1\cdots w_N)^{-1/N}$ 衡量模型拟合测试数据的程度，等价于预测下一个词时的加权平均分支因子，越低越好。}

\item 端到端语音识别中 CTC 与 Attention 的区别是什么？端到端相对传统混合系统有何优点？
\ansL{CTC 通过引入 blank 符号解决「一帧对应不到一个字」的对齐问题：输出是一串尖峰，某输出序列的概率等于填了 blank 之后所有可能路径的概率之和；它保留逐帧独立假设，可端到端训练。Attention（编码器--解码器）把输入序列编码为定长向量再由解码器解码，解码每一步动态关注输入的相关部分，不需要每帧独立的假设；但长序列给编码器带来巨大压力，噪声环境下通常不如 CTC，原始 LAS 模型还易出现对齐漂移。优点：端到端直接把声学特征映射到更大的单元（词、字）序列，而不是音素序列，省去了声学模型、词典、语言模型的分别训练与强制对齐，字级输出也便于下游 NLP 处理；混合式则需要声学模型＋词典＋语言模型。}

\end{qlist}


\chapstart{4}{语音合成}{选择题 12 · 填空题 10 · 判断题 6 · 简答题 5}{33}

\subsection{一、选择题}

\begin{qlist}

\item 语音合成（TTS）的任务是：
\optlisttwo{A. 把语音转写成文本}{B. 由文本自动生成语音}{C. 判断说话人身份}{D. 增强带噪语音}
\ans{B。语音合成指通过机械的、电子的方法产生人造语音的技术；TTS（Text-to-Speech）指自动把文本转换为声音。}

\item 语音合成的两个主要目标是：
\optlisttwo{A. 速度与体积}{B. 可懂度与自然度}{C. 信噪比与失真度}{D. 准确率与召回率}
\ans{B。可懂度指听者能听懂，自然度指听起来像不像真人说的。}

\item 关于主观评价指标，下列说法正确的是：
\optlist{A. MOS 取值范围为 $[-3,3]$}{B. CMOS 是绝对评分、范围 $1\sim 5$}{C. MOS 为 $1\sim 5$ 分制，CMOS 为相对评分、范围 $-3\sim 3$}{D. ABX 只能用于客观评价}
\ans{C。MOS（平均意见得分）分 5 级：优、良、一般、差、劣；CMOS（对比平均意见得分）为相对评分，$3$ 表示远优于、$0$ 无差别、$-3$ 远差于；ABX 测试播放 A、B 两段音频，由听者选 A、选 B 或认为无差别（X）。}

\item 下列属于语音合成客观评价指标的是：
\optlisttwo{A. MOS}{B. CMOS}{C. 梅尔倒谱失真（MCD）}{D. ABX}
\ans{C。客观指标是合成结果与真实录音（ground truth）之间的距离度量，包括 MCD、基频 $F_0$ 的 RMSE、时长的 RMSE。}

\item 波形拼接合成的缺点是：
\optlist{A. 无需大语料库}{B. 需要大规模精心设计的语音数据库，且不易控制韵律}{C. 合成质量不稳定}{D. 无法扩展}
\ans{B。其余缺点还包括：数据库中缺少类似音素样例时合成质量差、可扩展性差不易维护。这些正是统计参数合成要解决的问题。}

\item 统计参数语音合成（SPSS）的优点是：
\optlist{A. 需要大规模精心设计的语料库}{B. 无需大规模精心设计的语音数据库、合成质量稳定、可以控制韵律、可扩展性好}{C. 无法控制韵律}{D. 只能合成单字}
\ans{B。它先对语音做参数化表征，再用统计模型对声学参数建模，并以训练得到的模型构建合成系统，分基于 HMM 与基于 DNN 两条路线。}

\item 语音合成系统的正确分工是：
\optlist{A. 前端做声码器、后端做文本分析}{B. 前端做文本分析与韵律预测、后端做声学模型与声码器}{C. 前后端都做声学建模}{D. 前端输出波形、后端输出文本}
\ans{B。合成本质是把语言信息调制到声学信号上：前端提取语言学特征（音标与韵律信息的符号化表征），后端把语言学特征转换为声学特征、再转换为语音波形。有的研究把「文本分析＋韵律预测」合称语言分析。}

\item 韵律结构由小到大的三个层级是：
\optlist{A. 语调短语、韵律短语、韵律词}{B. 韵律词、韵律短语（中间短语）、语调短语}{C. 音节、词、句}{D. 声母、韵母、声调}
\ans{B。三个易混概念：语调是说话的腔调（一句话里快慢轻重的配置与变化）；声调又叫字调，普通话有阴平、阳平、上声、去声；重音起强调重点、突出情感的作用。}

\item 基于 HMM 的语音合成与语音识别的关系是：
\optlist{A. 完全无关}{B. 互为反向问题，训练都用 Baum-Welch、解码都用 Viterbi}{C. 合成用 Viterbi 训练、识别用 Baum-Welch}{D. 两者只能共用语言模型}
\ans{B。识别是 $\hat L=\arg\max_L P(O\mid L,\Lambda)$（由声音求文本），合成是 $\hat O=\arg\max_O P(O\mid L,\Lambda)$（由文本求声学参数）；训练算法相同（Baum-Welch），解码算法相同（基于动态规划的 Viterbi）。}

\item HMM 合成语音听起来平淡、音质不好的主要原因是：
\optlist{A. 训练数据不足}{B. 过平滑：状态参数分布是统计平均的结果，且决策树为浅层线性模型}{C. 声码器无法工作}{D. 基频提取错误}
\ans{B。HMM 假设各状态之间独立，状态参数分布是统计平均的结果；决策树作为浅层线性模型无法建模复杂映射关系，对输入特征与数据空间的线性分割会降低泛化能力。}

\item 基于 DNN 的语音合成中，时长模型与参数生成模型的输入分别是：
\optlist{A. 都是帧级输入}{B. 时长模型为音素级特征，参数生成模型为按音素时长展开后的帧级输入}{C. 都是音素级输入}{D. 时长模型为波形，参数生成模型为频谱}
\ans{B。一句话有多少音素就对应多少条时长模型的输入，有多少帧就对应多少条参数生成模型的输入。}

\item 端到端语音合成中，被称为「神经声码器」的是：
\optlisttwo{A. 前端的端到端模型}{B. 后端的端到端模型}{C. 文本分析模块}{D. 韵律预测模块}
\ans{B。前端的端到端模型把韵律预测、持续时长与语音特征模型统合到一个模型中，直接由上下文无关的符号预测语音特征；后端的端到端模型读取语音特征直接产生波形，即神经声码器。}

\end{qlist}

\subsection{二、填空题}

\begin{qlist}

\item \blank ：通过机械的、电子的方法产生人造语音的技术；\blank ：自动将文本转换为声音。
\ans{语音合成（Speech Synthesis）；TTS（Text-to-Speech）。}

\item 文本分析一般包含三步：\blank 、\blank 、\blank 。
\ans{分词与词性标注；文本归一化；字音转换（标音）。字音转换需处理协同发音、多音字消歧、儿化音标注、变调等。}

\item 文本归一化的目的是 \blank ，例如「前方 120M」应读作「前方一百二十米」。
\ans{把非标准词转换为发音明确的口语单词。方法包括基于规则转换、规则转换与深度学习结合。}

\item 韵律是 \blank 的接口；在语音 \blank 中通过 $F_0$、时长、能量等的变化实现，在语音 \blank 中反映节奏、强调、语速。
\ans{句法和发音；生成；感知。（语调调型隐含于基频模型——p53 蓝字。）}

\item 统计参数语音合成分 \blank 与 \blank 两条路线；SPSS 的四个优点是：无需大规模精心设计的语音数据库、\blank 、\blank 、可扩展性好易维护。
\ans{基于 HMM；基于 DNN；合成语音的质量稳定；可以控制韵律。}

\item \blank 语音合成兴起于 1990 年左右（规则驱动需大语料库），\blank 语音合成兴起于 1970 年左右。
\ans{波形拼接；参数。}

\item 如何解决波形拼接合成的问题？答案指向 \blank 语音合成。
\ans{统计参数（SPSS）。其优点：无需大规模精心设计的语音数据库、合成语音质量稳定、可以控制韵律、可扩展性好易于维护。}

\item 端到端语音合成直接读文本作为输入、不使用中间表示，直接输出 \blank ，其优化目标是 \blank 。
\ans{语音波形；整个系统（而不是各子系统），从而减少人工假设与多阶段建模造成的性能损失。}

\item VALL-E 把语音合成看作 \blank 任务而不是连续信号回归任务，仅需 \blank 秒的音频提示即可合成未见说话人的语音；其中第一个码本用 \blank 预测，其余码本用 \blank 预测。
\ans{条件语言模型；3；自回归（AR）；非自回归（NAR）。}

\item Tacotron 2 把前端的端到端模型 \blank 与后端的端到端模型 \blank 结合在一起。
\ans{Tacotron；WaveNet（的改良版）。}

\end{qlist}

\subsection{三、判断题}

\begin{qlist}

\item 语音合成与语音识别是互为反向的问题，但两者训练与解码所用的算法相同。
\ans{正确。训练都用 Baum-Welch（EM），解码都用基于动态规划的 Viterbi。}

\item MOS 是相对评分，CMOS 是绝对评分。
\ans{错误。正好相反：MOS 为绝对评分（$1\sim 5$ 级），CMOS 为相对评分（$-3\sim 3$）。}

\item 统计参数语音合成相比波形拼接更容易控制韵律，且不要求大规模精心设计的语音数据库。
\ans{正确。这是 SPSS 的主要优点，此外还有合成质量稳定、可扩展性好易于维护。}

\item 韵律只与音高有关，与时长、能量无关。
\ans{错误。韵律是超音段属性，从信号角度包括音高、音长和强度，在生成端通过 $F_0$、时长、能量共同实现。}

\item 声调与语调是同一概念，都指字调。
\ans{错误。声调（tone）又叫字调，普通话有阴平、阳平、上声、去声四个；语调（intonation）是说话的腔调，是一句话里快慢轻重的配置和变化，体现说话人的态度或口气。}

\item 端到端语音合成的优化目标是各个子系统分别最优。
\ans{错误。端到端的优化目标是整个系统，从而减少人工假设与多阶段建模导致的性能损失。}

\end{qlist}

\subsection{四、简答题}

\begin{qlist}

\item 画出（用文字描述）语音合成系统的整体框架，说明前端与后端各做什么。
\ansL{整体框架：文本 $\to$【前端：文本分析（分词与词性标注、文本归一化、字音转换）＋韵律分析（韵律结构预测、基频／时长／能量预测）】$\to$ 语言学特征（音标与韵律信息的符号化表征）$\to$【后端：声学模型把语言学特征转换为声学特征（谱参数、激励参数）＋声码器把声学参数还原为波形】$\to$ 语音。理论依据：语音是「通过语言调制、由人类发声器官产生的声学信号」，因此合成就是把语言信息调制到声学信号上。前端解决「怎么念」，后端解决「怎么发声」。}

\item 比较波形拼接与统计参数语音合成的基本思想与优缺点。
\ansL{波形拼接：先从自然语音中抽取词或句的语音基元并按规律组成语料库，合成时按文本查找合适基元，经拼接、韵律修饰得到语音。优点是音质自然（直接使用真人录音）；缺点是需要大规模精心设计的语音数据库、库中缺少类似音素样例时质量差、不易控制韵律、可扩展性差不易维护。统计参数合成（SPSS）：对输入语音做参数化表征，用统计模型对声学参数建模，再由模型构建合成系统。优点是无须大规模精心设计的语料库、合成质量稳定、可以控制韵律、可扩展性好易于维护；缺点是受模型精度与过平滑影响，音质不如波形拼接。二者代表了从「数据驱动拼接」到「模型驱动生成」的演进。}

\item 说明基于 HMM 的语音合成的系统框架（训练与合成两个阶段）。
\ansL{为什么用 HMM：语音产生本身是双重随机过程（气流经由的过程不可观测、发出的语音可观测），与 HMM 的结构同构。训练流程：①数据处理，提取文本特征与声学参数，训练数据为 〈文本特征，声学参数〉对；②模型训练，依次为单音素训练 $\to$ 上下文相关 HMM（三音素）$\to$ 决策树聚类 $\to$ 聚类后模型训练。合成流程：①文本处理得到输入文本特征；②基于似然最大化准则，用 HMM 生成声学参数；③用声码器把声学参数转换为最终语音。其固有缺陷是过平滑，导致合成语音平淡、音质不好。}

\item 什么是过平滑问题？它由哪些原因造成？
\ansL{过平滑指统计模型在参数生成时把细节抹平，使合成语音平淡、音质不好。原因有二：其一，HMM 基于建模假设各状态之间独立，状态参数分布是统计平均的结果，方差被平均掉，生成的参数轨迹过于平滑、丢失细节；其二，决策树作为浅层线性模型无法建模复杂的映射关系，对输入特征和数据空间的线性分割会降低模型的泛化能力。改进方向包括引入 DNN 等强非线性模型（DNN 合成）、以及端到端建模避免中间表征误差。}

\item 端到端语音合成有哪两类模型？各自的输入输出是什么？请举出代表性模型。
\ansL{端到端（E2E）TTS 直接读文本作为输入、不使用任何中间表示、直接输出语音波形，优化目标是整个系统。两类：①前端的端到端模型——把韵律预测模型、持续时长模型和语音特征模型统合到一个模型里，直接由上下文无关的符号预测语音特征；②后端的端到端模型——读取语音特征作为输入直接产生波形，又称神经声码器。代表模型：Tacotron 2（2018）把前端的端到端模型 Tacotron 与后端的端到端模型 WaveNet 改良版结合；FastSpeech（2019）用非自回归方式并利用文本与语音的单调对齐加速梅尔谱生成，解决自回归模型生成慢、鲁棒性差（跳词、重复）、可控性差的问题；VITS（2021）实现声学模型与声码器的统一建模，避免两段式框架中共用梅尔谱中间表征引入的误差；VALL-E 把合成看作条件语言模型任务。}

\end{qlist}


\chapstart{5}{语音增强}{选择题 12 · 填空题 10 · 判断题 6 · 简答题 5}{33}

\subsection{一、选择题}

\begin{qlist}

\item 单说话人场景下的受损语音模型可写为：
\optlisttwo{A. $y(n)=x(n)+h(n)$}{B. $y(n)=x(n)*h(n)+v(n)$}{C. $y(n)=x(n)\cdot v(n)$}{D. $y(n)=h(n)*v(n)$}
\ans{B。$x(n)$ 为干净语音、$h(n)$ 为房间脉冲响应（RIR）、$v(n)$ 为背景噪声；混响为乘性（卷积）干扰，噪声为加性干扰。}

\item 关于背景噪声与混响，下列说法正确的是：
\optlist{A. 背景噪声是乘性干扰、与干净语音相关性高}{B. 混响是加性干扰、相关性低}{C. 背景噪声是加性干扰、相关性更低；混响是乘性干扰、相关性更高}{D. 两者都是加性干扰}
\ans{C。混响本质是干净语音自身的多次延迟叠加，故与干净语音相关性更高；背景噪声来自其他声源，相关性更低。这一区别决定了谱减法只对加性噪声有效，而 IRM 等掩蔽目标不适用于混响。}

\item 狭义的语音增强与狭义的语音分离分别指：
\optlist{A. 去混响与去噪}{B. 语音去噪（语音与非语音分离）与说话人分离（多说话人分离）}{C. 语音识别与语音合成}{D. 单通道与多通道处理}
\ans{B。三类分离任务：语音去噪、语音去混响、说话人分离；前两类合称广义的语音增强。}

\item 基于传统信号处理的语音增强算法流程依赖的两个核心步骤是：
\optlist{A. 波束形成与后置滤波}{B. 静音检测（VAD）与噪声估计}{C. 特征提取与模型训练}{D. 编码与解码}
\ans{B。传统方法的目标就是算出一个抑制增益，而这个增益依赖于准确的噪声谱估计；通常利用 VAD 在静音帧更新噪声估计。只有噪声谱估得准，抑制增益才准。}

\item 谱减法中被广泛采用的假设是：
\optlist{A. 语音与噪声高度相关}{B. 语音与噪声不相关，功率谱可直接相减}{C. 噪声为零均值高斯}{D. 相位与幅度均可忽略}
\ans{B。由 $y(n)=x(n)+v(n)$ 且两者不相关得 $|Y(\omega)|^2\approx|X(\omega)|^2+|V(\omega)|^2$，于是 $|\hat X(\omega)|=\left(|Y(\omega)|^2-|\hat V(\omega)|^2\right)^{1/2}$，相位沿用带噪语音的相位。}

\item 谱减法最典型的副作用是：
\optlist{A. 计算量过大}{B. 产生「音乐噪声」（有节奏感、类似音乐声的残留）}{C. 无法实时实现}{D. 只能处理平稳噪声}
\ans{B。其缺点还包括适用信噪比范围窄，低信噪比时对语音可懂度损伤较大。广义谱减法用 $\max(|Y|^{2n}-\alpha|\hat V|^{2n},\ \beta|Y|^{2n})$ 避免出现负功率谱。}

\item 维纳滤波器的传输函数为：
\optlisttwo{A. $H(\omega)=\dfrac{P_v}{P_x+P_v}$}{B. $H(\omega)=\dfrac{P_x}{P_x+P_v}$}{C. $H(\omega)=\dfrac{P_x}{P_v}$}{D. $H(\omega)=1$}
\ans{B。其思想是寻找系数为实数的线性滤波器，使滤波后的信号与原干净语音的均方误差最小（故又称最小二乘／最小平方滤波器）。$P_x\gg P_v$ 时 $H\to1$ 保留，$P_x\ll P_v$ 时 $H\to0$ 压制。}

\item 频域语音增强的三个步骤是：
\optlist{A. 编码、传输、解码}{B. 特征提取、特征增强、波形重构}{C. 分帧、加窗、FFT}{D. 训练、验证、测试}
\ans{B。特征提取把时域信号转为频域特征（STFT 语谱图）；特征增强用神经网络增强频谱图的振幅；波形重构用增强后的振幅与原始相位经 ISTFT 还原时域波形。}

\item 基于映射（mapping）与基于掩蔽（mask）的框架，其网络输出分别是：
\optlist{A. 都为干净语音频谱}{B. 前者输出干净语音频谱，后者输出掩蔽值（如 $X/Y$）}{C. 前者输出掩蔽值，后者输出干净频谱}{D. 都输出掩蔽值}
\ans{B。映射把语音增强视为多变量回归问题，非线性回归函数由 DNN 参数化；掩蔽受听觉掩蔽现象与听觉场景分析启发，先估计掩蔽再与含噪频谱相乘得到干净语音估计。两者都明显优于谱减法与维纳滤波，且对非平稳噪声抑制良好。}

\item 关于两种经典掩蔽目标，说法正确的是：
\optlist{A. IBM 是软掩蔽、IRM 是硬掩蔽}{B. IBM 取 0 或 1（二值），常用交叉熵损失；IRM 是 $0\sim 1$ 的软比例，常用 MSE 损失}{C. 两者都用 MSE}{D. IRM 适用于混响}
\ans{B。$\mathrm{IBM}=1$（当 $\mathrm{SNR}(t,f)$ 大于阈值）否则为 0，是二分类问题，常用交叉熵；$\mathrm{IRM}=\left(\dfrac{|X|^2}{|X|^2+|V|^2}\right)^{\beta}$，$\beta$ 常取 0.5，常用 MSE。IRM 适用于加性噪声，不适用于混响（混响不是能量叠加）。}

\item 波束形成的基本思路是：
\optlist{A. 对单路信号做滤波}{B. 对各麦克风信号加权求和，通过调整权系数实现对目标方向的增强与对其他方向的抑制}{C. 对频谱做减法}{D. 用神经网络预测掩蔽}
\ans{B。加权系数决定增强效果；按权系数能否随信号统计特性自适应调整，分为固定波束形成与自适应波束形成。}

\item 延迟求和波束形成（DSB）包含的两个步骤是：
\optlisttwo{A. 编码与解码}{B. 延迟补偿与加权求和}{C. 训练与推理}{D. 分帧与加窗}
\ans{B。延迟补偿使各麦克风信号在时间上对齐，加权求和得到增强后的单通道输出；来自目标方向的声音同相相加被增强，其他方向因不同相被削弱。}

\end{qlist}

\subsection{二、填空题}

\begin{qlist}

\item 受损语音 $y(n)=x(n)*h(n)+v(n)$ 中，$h(n)$ 称为 \blank ，它由障碍物反射形成声音的持续传播。
\ans{房间脉冲响应（RIR，Room Impulse Response）。}

\item 语音增强的目的是去除受损语音中的 \blank 与 \blank ，得到目标干净语音；其意义一是提高语音 \blank ，二是提高 \blank 的性能。
\ans{背景噪声；混响；可懂度（清晰度）；语音识别等后端系统。注意不少增强技术在降噪的同时会引入失真，反而损伤可懂度。}

\item \blank ：用带噪信号的频谱减去噪声信号的频谱，得到 \blank 信号的频谱估计；其典型算法还有 \blank ，后者又被称为 \blank 滤波器（最小二乘／最小平方滤波器）。
\ans{谱减法；干净；维纳滤波；最小平方。}

\item \blank 是基于传统信号处理方法的典型算法；维纳滤波器也称 \blank 滤波器或 \blank 滤波器。
\ans{谱减法；最小二乘；最小平方。（传统方法最重要的是噪声估计，计算量小、通用性强。）}

\item 基于掩蔽的语音增强，增强阶段把 \blank 与 \blank 相乘，得到干净语音频谱的估计。
\ans{含噪语音的语谱图；模型输出的掩蔽值。}

\item 理想二值掩蔽 IBM 的二分类训练常用 \blank 损失；理想比率掩蔽 IRM 的估计常用 \blank 损失。
\ans{交叉熵（CE）；均方误差（MSE）。}

\item 麦克风阵列常见的两种几何形态是 \blank 与 \blank ；通过调整加权系数，对目标声源方向实现 \blank 、对其他方向实现 \blank 。
\ans{均匀线性阵列；均匀圆形阵列；增强；抑制。（此为波束形成 Beamforming。）}

\item TDOA 通过计算两路信号的 \blank 来估计，其峰值出现的延迟即为 $\hat\tau_s=\arg\max_\tau\phi_{12}(\tau)$。
\ans{互相关函数。因为 $\phi_{12}(\tau)=E[x_1(t)x_2(t+\tau)]=\phi_{11}(\tau-\tau_s)$，当 $\tau=\tau_s$ 时自相关取最大值。}

\item 复杂听觉场景下，人依靠 \blank 刺激驱动（被动）与 \blank 任务驱动（主动）的组合来选择特定的听觉物体。
\ans{自下而上；自上而下。}

\item 多目标训练指同一个网络同时输出干净语音 $X$ 与 \blank ；其技巧依据是映射与掩蔽在不同信噪比下 \blank 。
\ans{掩蔽（如 $X/Y$）；互补性很强。}

\end{qlist}

\subsection{三、判断题}

\begin{qlist}

\item 谱减法只需要正反傅里叶变换，实时实现较容易。
\ans{正确。这是其主要优点，但效果较差且易产生音乐噪声。}

\item 维纳滤波在均方误差最小意义下最优，因此不会残留噪声。
\ans{错误。它是均方误差最小意义下的最优线性滤波器，但实际中仍会有噪声残留（课件示例亦如此）。}

\item 深度学习语音增强擅长处理非平稳噪声，且不会有音乐噪声。
\ans{正确。深度学习具有非线性变换能力与强大的回归能力，增强后的语音更自然、没有音乐噪声。}

\item 基于映射的框架是把增强看作多变量回归问题，由 DNN 参数化非线性回归函数。
\ans{正确。这是课件对 mapping 框架的表述。}

\item IRM 可以直接用于去除混响。
\ans{错误。IRM 适用于加性噪声，混响是卷积而非能量叠加，因此不适用。}

\item 固定波束形成的权系数在设计好后保持不变，环境适应性较差。
\ans{正确。其计算简单、容易实现，但对噪声无自适应性；自适应波束形成可随信号统计特性自我调整。}

\end{qlist}

\subsection{四、简答题}

\begin{qlist}

\item 简述基于传统信号处理的语音增强算法流程，并说明为什么噪声估计是关键。
\ansL{流程：①分帧、加窗、FFT，得到每帧频谱 $Y(\omega)$；②用 VAD 判断该帧是否含语音；③只在静音帧更新噪声谱估计 $\hat V(\omega)$；④用 $Y$ 与 $\hat V$ 计算抑制增益（谱减法或维纳滤波）；⑤把增益乘回幅度谱，沿用含噪语音的相位，经 IFFT 与重叠相加得到增强语音。噪声估计之所以关键：传统方法的核心是算出一个抑制增益，而增益完全由噪声谱估计决定（谱减法为 $|Y|^2-|\hat V|^2$，维纳滤波为 $P_x/(P_x+P_v)$，其中 $P_x$ 也依赖噪声估计），噪声谱估得准增益才准。风险在于：VAD 一旦把语音帧误判为静音，就会把语音当成噪声学进去，之后反把那一段减掉。}

\item 比较谱减法与维纳滤波的基本原理与优缺点。
\ansL{谱减法：用带噪信号频谱减去噪声信号频谱得到干净信号频谱的估计：
\[|\hat X(\omega)|=\left(|Y(\omega)|^2-|\hat V(\omega)|^2\right)^{1/2}\]
相位沿用带噪语音（语音与噪声不相关的假设下 $|Y|^2\approx|X|^2+|V|^2$）。优点：简单，只需正反傅里叶变换，实时实现容易。缺点：适用信噪比范围窄，低信噪比时对可懂度损伤大；频谱直接相减会产生刺耳的「音乐噪声」。维纳滤波：寻找实数系数的线性滤波器使滤波后信号与原干净语音的均方误差最小，传输函数为
\[H(\omega)=\frac{P_x(\omega)}{P_x(\omega)+P_v(\omega)}\]
优点：性能优于谱减法，极大消减音乐噪声。缺点：原理较复杂、运算量相对较大；仍需估计噪声功率谱。两者共同的局限：不擅长处理非线性（混响等乘性、卷积）噪声与非平稳（时变）噪声，性能不佳、残留较多。}

\item 写出基于映射（mapping）的语音增强的训练与增强步骤。
\ansL{训练：①特征提取——用 STFT 提取含噪音频与干净音频的语谱图特征，含噪为训练的输入、干净为目标；②模型训练——用深度学习学习「含噪语谱图 $\to$ 干净语谱图」的非线性映射，训练非线性函数并保存模型。增强：①特征提取——用 STFT 提取含噪语音的语谱图；②特征增强——把含噪语谱图送入训练好的模型，得到干净音频特征的估计；③信号重构——把干净语谱图的估计值与原始相位作 ISTFT，重构出增强语音。相比传统方法，基于映射的框架效果有极大提升，对非平稳噪声有很好的抑制作用。}

\item 写出基于掩蔽（mask）的语音增强的训练与增强步骤，并说明它与映射框架的区别。
\ansL{训练：①特征提取——用 STFT 提取含噪与干净语谱图；②计算掩蔽值——计算从含噪音频到干净语音的掩蔽值作为标签；③模型参数更新——用深度学习学习「含噪语谱图 $\to$ 掩蔽值」的非线性映射并保存模型。增强：①特征提取——STFT 提取含噪语谱图；②计算掩蔽值——模型输出掩蔽值；③把含噪语音的语谱图与掩蔽值相乘，得到干净语音语谱图的估计；④用该估计与原始带噪语音的相位作 ISTFT，重构增强语音。与映射的区别：映射直接回归干净频谱（多变量回归），掩蔽把任务转为「该保留几成」的估计，网络输出的是掩蔽值，需再与含噪频谱相乘；掩蔽受听觉掩蔽现象与听觉场景分析启发，广泛用于降噪与多说话人分离。由于两者在不同信噪比下互补，可用多目标训练同时输出二者。}

\item 简述延迟求和波束形成的原理，并写出 TDOA 的估计方法。
\ansL{原理：对各个麦克风单元接收到的信号做加权求和，加权系数决定增强效果；通过调整加权系数，实现对目标声源方向的增强、对其他方向的抑制，使阵列对目标方向形成聚焦。DSB 包含两步：①延迟补偿——让每个麦克风的信号在时间上对齐（$D_i=D_0-(i-1)\hat\tau_s$）；②加权求和——输出 $y(t)=\sum_{i=1}^{M}w_ix_{si}(t-D_i)$，等权时 $w_i=1/M$。线性阵列中相邻麦克风路程差 $\xi=d\sin\theta_s$，故 $\tau_s=\xi/c=\dfrac{d\sin\theta_s}{c}$。TDOA 估计：利用互相关与自相关的关系 $\phi_{12}(\tau)=E[x_1(t)x_2(t+\tau)]=E[x_1(t)x_1(t+\tau-\tau_s)]=\phi_{11}(\tau-\tau_s)$，当 $\tau-\tau_s=0$ 时自相关取最大，故 $\hat\tau_s=\arg\max_\tau\phi_{12}(\tau)$，进而得到达方向 $\hat\theta_s=\sin^{-1}(c\hat\tau_s/d)$，并增强该方向到来的语音。}

\end{qlist}


\chapstart{6}{声纹识别}{选择题 12 · 填空题 8 · 判断题 6 · 简答题 4}{30}

\subsection{一、选择题}

\begin{qlist}

\item 声纹识别的任务是：
\optlist{A. 把语音转写成文本}{B. 根据说话人的声音特性判断其身份}{C. 判断语音所属语种}{D. 由文本生成语音}
\ans{B。声纹识别（说话人识别）是生物识别（指纹、掌纹、人脸、虹膜等）的一种：输入语音，输出说话人身份。}

\item 下列任务与输出信息的对应关系错误的是：
\optlisttwo{A. 语音识别 $\to$ 文本信息}{B. 情感识别 $\to$ 情感信息}{C. 语种识别 $\to$ 说话人身份}{D. 声纹识别 $\to$ 说话人身份信息}
\ans{C。语种识别输出的是语音所属语种，不是说话人身份。}

\item 与说话人身份相关的语音属性中，低级属性与高级属性分别指：
\optlist{A. 韵律与口音、声道频谱与共振峰轨迹}{B. 与发声器官的生理方面相关（声道频谱、共振峰轨迹）、与说话方式的行为差异相关（韵律、口音、用词）}{C. 音高与音长}{D. 文本相关与文本无关}
\ans{B。}

\item 说话人辨认（Identification）与说话人确认（Verification）的对比形式分别是：
\optlisttwo{A. 1:1 与 1:N}{B. 1:N 与 1:1}{C. 都是 1:N}{D. 都是 1:1}
\ans{B。辨认判定待识别说话人来自注册集中的哪一个（$1:N$）；确认判定待识别语音是否来自注册集中的某个说话人（$1:1$，成对）。}

\item 关于 FRR、FAR、EER，下列说法正确的是：
\optlist{A. FRR 是误接受率、FAR 是误拒绝率}{B. EER 是 FAR 与 FRR 相等时的错误率，越小越好}{C. EER 越大越好}{D. FAR 与 FRR 相等时系统不可用}
\ans{B。FRR（误拒绝率）是被错误拒绝的有效身份声明的比例（本该匹配成功却判失败）；FAR（误接受率）是错误接受冒名顶替者声明的比例（本该失败却判成功）；调整阈值使 $\mathrm{FAR}=\mathrm{FRR}$ 时的错误率即 EER。阈值变化时两者的权衡画成 DET 曲线。}

\item 声纹识别系统基本框架的三个部分是：
\optlisttwo{A. 特征提取、注册、测试}{B. 训练、验证、测试}{C. 前端、后端、解码}{D. 编码、传输、解码}
\ans{A。特征提取把原始语音转换为以说话人信息为主的特征向量序列；注册用注册语音建立说话人模型；测试从待测语音提取特征并与各注册说话者比较判决。}

\item 通用背景模型（UBM）的主要作用是：
\optlist{A. 直接作为说话人模型}{B. 用大量说话人无关数据训练的先验模型，解决单人语音太少、训练不出可靠 GMM 的问题}{C. 用于语种识别}{D. 用于特征降维}
\ans{B。实际场景中每个说话人的语音数据很少，且训练与测试存在信道失配；UBM 用 EM 算法在说话人无关数据上训练，再通过 MAP 等自适应算法得到目标说话人模型。}

\item 使用 MAP 自适应的好处不包括：
\optlist{A. 减少注册新说话人所需的语音时长／数量}{B. 说话人相关 GMM 与 UBM 模型复杂度相同，可计算更可靠的对数似然分数}{C. 可以使用高斯选择从而更快}{D. 可以完全不需要任何注册语音}
\ans{D。MAP 自适应仍需少量注册语音，只是大幅减少所需时长与数量。}

\item GMM 超向量的构造方法是：
\optlist{A. 把所有高斯分量的均值向量按固定顺序首尾相连}{B. 把协方差矩阵求逆}{C. 把权重求和}{D. 把多个说话人的模型平均}
\ans{A。做法是假设说话人与 UBM 的权重和协方差相同、只调整均值向量，再把各分量均值串成一条超长向量，作为说话人的定长表征。
}

\item x-vector 包含的四个函数块依次是：
\optlist{A. 特征提取、编码、解码、打分}{B. 帧处理（TDNN）、池化、嵌入提取、说话人分类}{C. 训练、注册、测试、判决}{D. 分帧、加窗、FFT、DCT}
\ans{B。帧处理器用多层网络（如 TDNN）提取深度特征；池化层把变长序列压成固定维向量；嵌入提取层中一层为瓶颈层，瓶颈层的输出（非线性层之前）就是 x-vector；最后用判别性损失按说话人标签分类。}

\item 时间池化（Temporal pooling）通常计算的是：
\optlisttwo{A. 最大值}{B. 特征向量各维的均值和标准差}{C. 中位数}{D. 一阶差分}
\ans{B。时间池化把特征向量序列转换为单个固定维向量，最常见做法是取均值与标准差。}

\item x-vector 中真正作为说话人嵌入输出的是：
\optlist{A. 池化层的输出}{B. 瓶颈层的输出（非线性激活之前）}{C. softmax 层的输出}{D. 输入特征}
\ans{B。嵌入提取由若干相连网络层组成，其中一层设计为瓶颈层，瓶颈层的输出（非线性层之前）即 x-vector 说话人嵌入，是固定长度向量，捕捉说话人的长期声纹特征。}

\end{qlist}

\subsection{二、填空题}

\begin{qlist}

\item 按输入模式，声纹识别分为 \blank 与 \blank ；前者性能更高，后者便利性更高。
\ans{文本相关（Text-dependent）；文本无关（Text-independent）。文本相关时注册与测试话语共享（至少部分）相同内容；文本无关时可用任意文本甚至不同语言。}

\item 与发声器官的 \blank 方面相关的属性称为低级属性（声道频谱、共振峰轨迹）；与说话方式的 \blank 差异相关的属性称为高级属性（韵律、口音、话语中词的选择）。
\ans{生理；行为。}

\item 说话人确认的任务是确定未知说话人是否是 \blank 的特定人。
\ans{一组人中（注册集中）。}

\item GMM-UBM 范式中，说话人相关的 GMM 由 \blank 唯一表示，方法是仅调整 \blank 向量，而 \blank 和 \blank 保持与 UBM 相同。
\ans{均值向量；均值；权重；协方差矩阵。（由此把各分量均值串接成 GMM 超向量。）}

\item 声纹识别系统：输入 \blank ，输出 \blank ；与 \blank （输出转写文本）形成对比。
\ans{语音；说话人身份；语音转写。}

\item 声纹识别与 \blank 的区别：前者输出语音携带的说话人身份信息，后者输出语音转写文本。
\ans{语音识别。}

\item 说话人 \blank （Speaker Identification）确定未知说话人是否是一组人中的特定人；说话人 \blank （Speaker Verification）判断是否来自注册集中的某个说话人。
\ans{辨认（Identification）；确认（Verification）。}

\item \blank 的说话人识别（Text-dependent）：注册与测试话语共享（至少部分）相同内容，性能更高；\blank 的说话人识别（Text-independent）：可用任意文本，便利性更高。
\ans{文本相关；文本无关。}

\end{qlist}

\subsection{三、判断题}

\begin{qlist}

\item 声纹识别属于生物识别，其输出是说话人身份。
\ans{正确。与语音识别（输出文本）、情感识别（输出情感）、语种识别（输出语种）不同。}

\item 文本无关的说话人识别性能一定高于文本相关方式。
\ans{错误。文本相关性能更高（注册与测试内容相同或部分相同），文本无关便利性更高。}

\item EER 越小，说明系统性能越好。
\ans{正确。EER 是调整阈值使 FAR 与 FRR 相等时的错误率，越小越好。}

\item UBM 是用目标说话人的数据训练出来的模型。
\ans{错误。UBM 用大量说话人无关的数据通过 EM 算法训练，是说话人无关的先验模型；目标说话人模型由 UBM 经 MAP 自适应得到。}

\item GMM 超向量通过同时调整权重、均值与协方差得到。
\ans{错误。为简化计算并减少数据需求，假设权重和协方差与 UBM 相同，只调整均值向量，再把各分量均值串接成超向量。}

\item x-vector 的池化层把变长的特征序列转换为一个固定维向量。
\ans{正确。时间池化通常取各维的均值与标准差。}

\end{qlist}

\subsection{四、简答题}

\begin{qlist}

\item 说明声纹识别系统的基本框架，并比较说话人辨认与说话人确认。
\ansL{系统三部分：①特征提取——把原始语音转换为以说话人信息为主的特征向量序列（目的为保留说话人信息、去除噪声与其他干扰）；②注册——用注册语音建立说话人模型，供与测试段对照；③测试——从待测语音提取特征，与每个注册说话者比较并作出判决。辨认（Identification）：判定待识别说话人来自注册集中的哪一个，形式为 $1:N$，取打分最高的说话人模型。确认（Verification）：判定待识别语音是否来自注册集中的某个说话人，形式为 $1:1$（成对），判决基于对数似然比与阈值比较。}

\item 为什么需要通用背景模型（UBM）？MAP 自适应解决了什么问题？
\ansL{直接为每个说话人训练 GMM 存在两个问题：①实际场景中每个说话人的语音数据很少，无法训练出高效可靠的 GMM；②复杂声学环境与信道下，训练语音与测试语音存在失配。解决思路：先用 EM 算法在大量说话人无关的数据上训练一个说话人无关的领域模型／通用背景模型（UBM）作为先验；再以 UBM 为初始模型，用最大后验概率（MAP）自适应得到目标说话人模型。MAP 自适应可大幅减少注册新说话人所需的语音时长与数量；由于说话人相关 GMM 与 UBM 具有相同的模型复杂度，可以计算更可靠的对数似然分数；还可使用高斯选择加快计算。自适应时权重、均值向量和协方差矩阵都可调整：均值适应使高斯分量偏移，协方差适应改变分量的大小（不确定性）。}

\item 写出 GMM-UBM 的似然比打分公式，并解释各符号含义。
\ansL{给定测试段与说话人模型 $\theta$ 的平均对数似然为 $\mathcal L(Y\mid\theta)=\dfrac{1}{T}\sum_{t=1}^{T}\log p(x_t\mid\theta)$。判决用对数似然比（说话人模型 vs UBM）：$\Lambda(Y)=\mathcal L(Y\mid\theta_{\text{speaker}})-\mathcal L(Y\mid\theta_{\text{UBM}})\underset{H_1}{\overset{H_0}{\gtrless}}\eta$。其中 $Y$ 为测试段特征序列、$T$ 为帧数、$x_t$ 为第 $t$ 帧特征、$\theta_{\text{speaker}}$ 为目标说话人的 GMM 模型（由 UBM 经 MAP 自适应得到）、$\theta_{\text{UBM}}$ 为通用背景模型、$\eta$ 为判决阈值。$\Lambda(Y)$ 大于阈值则接受 $H_0$（判为同一说话人），否则拒绝。其思想是：比起「大众平均」模型，这段话在多大程度上更像「他」。}

\item 介绍 x-vector 的四个函数块及其作用。
\ansL{①帧处理器：使用多层神经网络（如延时神经网络 TDNN）实现帧级处理，从声学特征（MFCC 或对数 mel-filterbank 系数）中提取深度说话人特征；TDNN 提取上下文的能力取决于感受野宽度，可用空洞卷积或堆叠更多空洞一维卷积层扩大视野。②池化层：把变长的特征向量序列转换为一个固定维向量，最常用时间池化，即对每维特征计算均值与标准差。③嵌入提取：由若干相互连接的网络层组成，其中一层设计为瓶颈层，瓶颈层的输出（非线性层之前）就是 x-vector 说话人嵌入，为固定长度向量，捕捉说话人的长期声纹特征。④说话人分类：输出层每个节点代表一个说话人，用 softmax 归一化，用交叉熵损失衡量输出与说话人标签的差异，以判别性方式训练整个网络。x-vector 的目的是用监督方式训练的深度神经网络导出句子级别的嵌入表示。}

\end{qlist}


\end{document}
